\section{Problem Setup}

\subsection{Top-1 routing and expert usage}

Let $h_i\in\R^d$ be the route-relevant representation of example $i$. The
top-1 router has one prototype $u_e\in\R^d$ for expert
$e\in\{1,\ldots,E\}$. It computes
\begin{equation}
 s_e(h_i)=\langle h_i,u_e\rangle,
 \qquad
 a(h_i)=\argmax_e s_e(h_i).
\end{equation}
The empirical expert usage is
\begin{equation}
 \widehat p_e=\frac{1}{N}\sum_{i=1}^N \mathbf{1}\{a(h_i)=e\}.
\end{equation}
The count-level target is $\widehat p_e\approx 1/E$ for all experts. This paper
asks when that target arises from the top-1 geometry itself, not from an
external capacity constraint.

\subsection{Four notions of uniformity}

\begin{definition}[Feature-frequency uniformity]
Feature labels $f\in\{1,\ldots,F\}$ are frequency-uniform if
$\Prob(f)=1/F$ for every $f$.
\end{definition}

\begin{definition}[Directional uniformity]
After centering and normalizing route-relevant states,
\begin{equation}
 z_i=\frac{h_i-\mu}{\|h_i-\mu\|_2},
\end{equation}
the token-side representation is directionally uniform if $z_i$ is
approximately uniform on $\Sphere^{d-1}$.
\end{definition}

\begin{definition}[Hidden spectral uniformity]
Let $\bar h_i=h_i-\mu$ and
\begin{equation}
 C=\frac{1}{N}\sum_i \bar h_i\bar h_i^\top.
\end{equation}
The representation is spectrally uniform when $C$ has no dominant
eigendirection, ideally $C=\sigma^2 I$. We use covariance condition number and
effective rank as diagnostics.
\end{definition}

\begin{definition}[Expert-usage uniformity]
The router is usage-uniform when $\widehat p_e\approx 1/E$ for all experts.
\end{definition}

The order matters. Frequency uniformity is a data property. Directional
uniformity is the token-side condition directly seen by the inner-product
router. Hidden spectral uniformity is a second-order diagnostic that can help
detect or repair anisotropy, but it is not the same as directional uniformity.
Usage uniformity is the router outcome.

\section{Theory}

\subsection{A sufficient condition for natural top-1 balance}

Normalize token features and expert prototypes to the sphere:
$z\in\Sphere^{d-1}$ and $u_e\in\Sphere^{d-1}$. Prototype $u_e$ owns the
spherical top-1 cell
\begin{equation}
 V_e=\{z\in\Sphere^{d-1}:\langle z,u_e\rangle\ge
 \langle z,u_j\rangle,\;\forall j\}.
\end{equation}
Tie boundaries are measure-zero under the continuous spherical uniform
distribution, so the tie-breaking rule does not affect the result below.

\begin{proposition}[Equal-measure cells are sufficient]
\label{prop:equal_measure_cells}
Assume $z\sim\mathrm{Unif}(\Sphere^{d-1})$ and the cells $V_1,\ldots,V_E$ form
an equal-measure partition of the sphere up to measure-zero boundaries. Then
\begin{equation}
 \Prob(a(z)=e)=\frac{1}{E},\qquad \forall e,
\end{equation}
and for independent examples, $\E[\widehat p_e]=1/E$.
\end{proposition}

\begin{proof}
The event $a(z)=e$ is exactly $z\in V_e$. The uniform distribution on the sphere
assigns probability according to spherical measure. Since the $E$ cells cover
the sphere up to measure-zero boundaries and have equal measure, each cell has
probability mass $1/E$.
\end{proof}

\begin{corollary}[Regular-simplex construction]
If $E\le d+1$ and $u_1,\ldots,u_E$ are the vertices of a regular simplex
centered at the origin, then the induced top-1 cells have equal spherical
measure. Therefore uniform spherical features have uniform expected top-1
expert usage.
\end{corollary}

\begin{proof}
A regular simplex is invariant under all permutations of its vertices. These
permutations are realized by orthogonal transformations of the span of the
simplex and can be extended by the identity on the orthogonal complement in
$\R^d$. Thus every cell can be mapped to every other cell by a measure-preserving
orthogonal map. The cells are congruent and have equal spherical measure.
\end{proof}

\subsection{Finite tokens do not fix prototype-cell imbalance}

For any fixed prototype geometry, define the population cell mass
$p_e=\Prob(z\in V_e)$. Then
\begin{equation}
 \widehat p_e-\frac{1}{E}
 =
 \left(p_e-\frac{1}{E}\right)
 +
 \left(\widehat p_e-p_e\right).
\end{equation}
The first term is the prototype-cell geometry error. The second term is
finite-token sampling noise. Increasing $N$ reduces the second term but does not
change the first.
